% Subsecção da Secção "Análise de Dados" (02_analise_dados.tex).
% Baseado em docs/02_preprocessamento/ do repositório.

\subsection{Pré-processamento}
\label{sec:preprocessamento}

O pré-processamento seguiu a ordem imposta pelo desenho da Secção~\ref{sec:formulacao}:
primeiro, apenas regras que não dependem de nenhuma estatística dos dados (remoção de
linhas estruturais e partição temporal); depois, todas as análises de qualidade e decisões
de limpeza feitas \textbf{só no período de desenvolvimento}; por fim, as transformações
aprendidas, ajustadas apenas no treino de cada \textit{fold}. Distinguem-se dois tipos de
operação: \textbf{regras determinísticas} (conversão de formatos, códigos que significam
``em falta'', exclusões por regra fixa), que não aprendem nada com os dados e são aplicadas
da mesma forma ao teste; e \textbf{transformações aprendidas} (imputação, \textit{clipping},
padronização, \textit{encoding}), que ficam num \textit{pipeline} e nunca veem dados de
validação ou de teste no ajuste. Todo o código está em classes do pacote
\texttt{src/lending/}, usadas pelos \textit{notebooks}, e cada exclusão é registada no
momento em que é feita, de forma que as tabelas de exclusões são geradas e não escritas à
mão.

% ---------------------------------------------------------------------
\subsubsection{Ficheiros de origem, schema e unidades}

Os dois ficheiros foram convertidos para \textit{pickle} com todas as colunas lidas como
texto, para que nenhum tipo fosse adivinhado pelo \textit{pandas}
(Tabela~\ref{tab:origem}). O tipo de cada uma das 151 colunas do ficheiro de aceites foi
declarado à mão, a partir do formato dos valores e do dicionário de dados: 112 numéricas,
17 categóricas, 11 datas no formato ``Mon-AAAA'', 5 de texto livre, 3 indicadores
\textit{Y/N}, o identificador, o prazo (``36 months'') e os anos de emprego (``< 1 year''
a ``10+ years'', convertidos em 0 a 10). A conversão é feita por um \textit{parser} que
dá erro perante um valor com formato inesperado, em vez de o transformar silenciosamente
em valor em falta; no desenvolvimento, nenhum valor ficou fora do formato. As unidades
foram confirmadas pelos intervalos observados: as percentagens estão na escala 0--100
(\texttt{int\_rate} entre 5 e 26), e \texttt{fico\_range\_high} é sempre
\texttt{fico\_range\_low} mais 4 pontos (mais 5 no escalão 845--850).

\begin{table}[!t]
\caption{Ficheiros de origem}
\label{tab:origem}
\centering
\footnotesize
\begin{tabular}{|l|r|r|l|}
\hline
\textbf{Ficheiro} & \textbf{Linhas} & \textbf{Col.} & \textbf{Período} \\
\hline
\textit{accepted} & 2\,260\,701 & 151 & \texttt{issue\_d} 2007-06 a 2018-12 \\
\hline
\textit{rejected} & 27\,648\,741 & 9 & pedido 2007-05 a 2018-12 \\
\hline
\end{tabular}
\end{table}

As datas mais recentes do ficheiro de aceites mostram que a extração foi feita por volta
de março de 2019 (Secção~\ref{sec:split}). O ficheiro não está por ordem cronológica: é a
junção de 16 exportações da LendingClub, e começa em 2015.

% ---------------------------------------------------------------------
\subsubsection{Linhas estruturais e partição temporal}

O ficheiro de aceites tem 33 linhas que não são empréstimos: as duas linhas de totais no
fim de cada exportação (``\textit{Total amount funded in policy code\ldots}'') e um
separador (``\textit{Loans that do not meet the credit policy}''), espalhadas pelo meio do
ficheiro. São identificadas por uma regra estrutural (identificador não numérico e todas
as outras colunas vazias) e removidas antes da partição. O identificador \texttt{id} é
único nas restantes 2\,260\,668 linhas. A partição usa os cortes de calendário da
Secção~\ref{sec:split}: 466\,345 empréstimos no desenvolvimento (2007-06 a 2014-12),
421\,095 no teste final (2015) e 1\,373\,228 no período reservado (2016--2018). O
desenvolvimento e o teste ficam guardados à parte, ainda em texto, e o teste só volta a
ser aberto para receber as mesmas regras de limpeza.

% ---------------------------------------------------------------------
\subsubsection{Valores em falta}

A análise procurou o \textbf{significado} de cada valor em falta, e não só a sua
frequência. Distinguem-se cinco casos: ausência de evento (ex.:
\texttt{mths\_since\_last\_delinq} em falta significa que nunca houve atraso), ausência de
denominador (ex.: \texttt{revol\_util} sem crédito rotativo), campos que só existem em
alguns empréstimos (\textit{hardship}, \textit{settlement}), campos só dos pedidos
conjuntos e campos não registados.

Foram encontrados dois problemas. Primeiro, em 2007--2008, ``nunca aconteceu'' estava
guardado como 0 em \texttt{mths\_since\_last\_record} e \texttt{mths\_since\_last\_delinq}:
``0 meses desde o último registo'' é impossível com 0 registos, pelo que 1\,275 e 820
valores passaram a valor em falta. Segundo, \textbf{37 colunas do \textit{bureau}} só
começaram a ser registadas em 2012 e estão completas a partir de 2013; estão vazias em todo
o treino da classificação (2007--2011) e incluem oito das variáveis sugeridas no enunciado.
Esta constatação levou às decisões descritas na Secção~\ref{sec:split}: a classificação não
as usa e a regressão é treinada em 2013--2014, onde faltam em menos de 1\% das linhas.
Além disso, 31 colunas estão totalmente vazias no desenvolvimento, entre elas
\texttt{member\_id} e as do pedido conjunto, porque não há pedidos conjuntos antes de 2015.

% ---------------------------------------------------------------------
\subsubsection{Duplicados e mutuários repetidos}

Não há empréstimos duplicados: o \texttt{id} é único, o \texttt{url} é coerente com ele e
não há linhas repetidas. Os mutuários repetidos, pelo contrário, não podem ser
identificados, porque \texttt{member\_id} está vazio em todo o ficheiro. Uma chave
aproximada (data da primeira linha de crédito, código postal, estado, rendimento, emprego e
situação habitacional) estima que 2\,106 empréstimos (0,48\%), em 1\,049 grupos, sejam de
pessoas com mais de um empréstimo, com um intervalo mediano de oito meses. Não se remove
nenhum, por serem empréstimos reais; a consequência --- o mesmo mutuário pode estar no
treino e na validação --- é uma limitação reconhecida.

% ---------------------------------------------------------------------
\subsubsection{Valores suspeitos e consistência}

Foram verificadas 22 regras de intervalo e de consistência entre colunas. O critério de
decisão foi: corrigir apenas quando o valor correto é conhecido ou o valor é um código;
excluir apenas quando o registo contradiz o que o enunciado pede; nos restantes casos,
manter e documentar. Daqui resultaram três ações:
\begin{itemize}
  \item \textbf{221 empréstimos com condições modificadas} foram excluídos: 139 têm
  \texttt{int\_rate} = 6,00\% em graus B a G, abaixo da taxa mínima de qualquer desses
  graus (ex.: um D4 a 6\% quando os D4 do mês tinham 17,14\%), e 203 têm uma prestação que
  difere em mais de 1~USD da fórmula de amortização. O enunciado pede os valores
  \textbf{originais} de \texttt{int\_rate} e \texttt{installment}, e \texttt{int\_rate} é o
  alvo da regressão.
  \item O valor 9\,999\,999 em \texttt{tot\_hi\_cred\_lim} e \texttt{total\_rev\_hi\_lim} é
  o teto do campo e não um limite real (3 e 2 valores), pelo que passou a valor em falta.
  \item Os valores extremos plausíveis (\texttt{revol\_util} até 892\%, coerente com o saldo
  e o limite; rendimentos de vários milhões) foram mantidos e são tratados no
  \textit{pipeline}. Pequenas incoerências entre agregados do \textit{bureau} foram
  documentadas, sem correção, por não se saber qual das colunas está certa.
\end{itemize}

% ---------------------------------------------------------------------
\subsubsection{Coortes e tabela de exclusões}

A limpeza é comum a todas as tarefas e produz um conjunto de desenvolvimento com
466\,124 empréstimos. A partir dele, cada tarefa aplica a sua coorte
(Tabela~\ref{tab:exclusoes}). Na regressão, os empréstimos de 2007-06 a 2012-12 ficam fora
do treino, pela razão descrita na Secção~\ref{sec:split}; restam 370\,278 (134\,767 de 2013
e 235\,511 de 2014). Na classificação, os empréstimos de 60 meses e o embargo de 2012--2014
deixam 31\,522 empréstimos de 36 meses, todos com o prazo terminado: 27\,646
\textit{Good standing} (87,7\%) e 3\,876 \textit{Default} (12,3\%). Os passos ``desfecho
desconhecido'' removem 0 empréstimos no treino, porque todos os \textit{Current} são de 60
meses. Entre os empréstimos da classificação há 2\,437 ``\textit{Does not meet the credit
policy}'', que entram pelo seu desfecho e cuja indicação nunca é usada como preditor. O
clustering de mutuários usa todo o conjunto limpo.

\begin{table}[!t]
\caption{Tabela de exclusões (ficheiro accepted)}
\label{tab:exclusoes}
\centering
\footnotesize
\begin{tabular}{|p{0.56\columnwidth}|r|r|}
\hline
\textbf{Passo} & \textbf{Removidas} & \textbf{Ficam} \\
\hline
Ficheiro original & -- & 2\,260\,701 \\
\hline
Linhas que não são empréstimos & 33 & 2\,260\,668 \\
\hline
Teste final, 2015 (partição) & 421\,095 & 1\,839\,573 \\
\hline
Reservado, 2016--2018 (partição) & 1\,373\,228 & 466\,345 \\
\hline
Condições modificadas & 221 & \textbf{466\,124} \\
\hline
\multicolumn{3}{|l|}{\textit{Regressão}} \\
\hline
Fora de 2013-01 a 2014-12 (partição) & 95\,846 & \textbf{370\,278} \\
\hline
\multicolumn{3}{|l|}{\textit{Classificação (treino)}} \\
\hline
Prazo de 60 meses & 128\,228 & 337\,896 \\
\hline
Embargo 2012--2014 (partição) & 306\,374 & 31\,522 \\
\hline
Desfecho desconhecido (\textit{Current}) & 0 & \textbf{31\,522} \\
\hline
\multicolumn{3}{|l|}{\textit{Teste final (mesmas regras)}} \\
\hline
Empréstimos de 2015 & -- & 421\,095 \\
\hline
Condições modificadas & 207 & \textbf{420\,888} \\
\hline
\end{tabular}
\end{table}

% ---------------------------------------------------------------------
\subsubsection{Teste final}

O teste final recebeu exatamente as mesmas regras de limpeza, pela mesma ordem e sem
qualquer ajuste: conversão de tipos, zeros que significam ``nunca aconteceu'', exclusão
das condições modificadas e códigos de teto. Para garantir que são as mesmas, as regras
foram reunidas numa única classe, aplicada aos dois conjuntos. Antes de a usar no teste,
verificou-se que, aplicada ao desenvolvimento em texto, a classe reproduz o conjunto limpo
célula a célula. Do teste só foram registadas as contagens exigidas pela tabela de
exclusões (Tabela~\ref{tab:exclusoes}); nenhuma estatística ou distribuição foi observada.
As coortes do teste só são construídas na avaliação final, com as mesmas regras; aí, os
eventuais empréstimos \textit{At risk} contam como \textit{Default}
(Secção~\ref{sec:split}).

% ---------------------------------------------------------------------
\subsubsection{Pedidos rejeitados}

O ficheiro de rejeitados não tem linhas estruturais. Foi partido com os mesmos cortes, pela
data do pedido: 3\,450\,084 pedidos no desenvolvimento, depois de afastados 49 pedidos de
maio de 2007 (anteriores ao primeiro empréstimo aceite), 2\,859\,379 de 2015 e
21\,339\,229 de 2016--2018. As colunas foram convertidas pelo mesmo \textit{parser} e
renomeadas para o nome da coluna equivalente dos aceites (ex.: \textit{Debt-To-Income
Ratio}~$\rightarrow$~\texttt{dti}, de ``38.64\%'' para 38,64), para que as comparações entre
os dois ficheiros sejam diretas. A Tabela~\ref{tab:rejeitados} resume as regras de
limpeza; o resultado são 3\,424\,228 pedidos.

\begin{table}[!t]
\caption{Limpeza dos pedidos rejeitados (desenvolvimento)}
\label{tab:rejeitados}
\centering
\footnotesize
\begin{tabular}{|p{0.38\columnwidth}|r|p{0.32\columnwidth}|}
\hline
\textbf{Regra} & \textbf{Linhas} & \textbf{Razão} \\
\hline
Excluir \texttt{policy\_code} = 2 & 6\,241 & Produtos não públicos; os aceites são todos 1 \\
\hline
Excluir linhas iguais em todas as colunas & 19\,615 & Mesmo pedido submetido várias vezes \\
\hline
\texttt{risk\_score} = 0 $\rightarrow$ em falta & 86\,038 & Indicado no enunciado \\
\hline
\texttt{dti} = $-1$, 9\,999, 99\,999, 199\,998 $\rightarrow$ em falta & 264\,300 & Códigos, não rácios \\
\hline
\texttt{dti} $>$ 1000\% $\rightarrow$ em falta & 8\,253 & Rendimento mal introduzido (ex.: 4042,11\%) \\
\hline
Montante = 0 $\rightarrow$ em falta & 2 & Sem montante \\
\hline
\end{tabular}
\end{table}

O DTI acima de 1000\% significaria dívidas mensais superiores a dez vezes o rendimento; os
valores entre 100\% e 1000\% (cerca de 72\,000) são plausíveis num pedido rejeitado e foram
mantidos, pelo que as análises do DTI dos rejeitados devem usar estatísticas robustas. O
achado mais importante é uma \textbf{mudança de escala do \textit{risk score}} a
2013-11-06: até essa data, os valores vão de cerca de 420 a 850, como o FICO; a partir dela,
o mínimo passa a 501 e o máximo chega a 990, o intervalo do VantageScore. A mudança
acontece num dia exato, pelo que reflete uma alteração da pontuação registada e não dos
requerentes. Como não há conversão exata entre as escalas, os valores não foram alterados e
foi acrescentada a coluna \texttt{risk\_score\_scale}; a comparação com o FICO dos aceites só
é válida no período FICO (1\,402\,189 pedidos). Os anos de emprego também não são fiáveis
a partir de 2010: a proporção de ``< 1 year'' passa de cerca de 27\% para 85\%, o que
sugere um valor por omissão do formulário. Os pedidos acima do limite da plataforma
(35\,000~USD; 17\,711 pedidos) foram mantidos, por serem pedidos reais.

% ---------------------------------------------------------------------
\subsubsection{\textit{Pipeline} de pré-processamento}

As variáveis de cada tarefa são lidas diretamente da tabela de elegibilidade da Secção~\ref{sec:formulacao}:
entram as elegíveis, exceto as condicionais (geografia, por razões de \textit{fairness}), o
texto livre e as datas, que só entrarão transformadas, como novas variáveis. Uma verificação
automática dá erro se alguma coluna proibida pelo enunciado aparecer entre as de uma tarefa.
Ficam 90 variáveis na regressão, 98 na classificação e 73 no clustering.

As transformações aprendidas estão num \textit{pipeline} do \textit{scikit-learn}, com três
ramos: (i) montantes em USD: logaritmo, \textit{clipping} aos percentis 0,5\% e 99,5\%,
imputação pela mediana com indicador de ``em falta'' e padronização; (ii) restantes
numéricas: os mesmos passos sem logaritmo; (iii) categorias: ``em falta'' como categoria
própria e \textit{one-hot}, com as categorias abaixo de 0,5\% do treino agrupadas. Em cada
ramo, as colunas vazias no treino são retiradas (ex.: as 37 colunas do \textit{bureau} no
treino da classificação). O indicador de ``em falta'' distingue ``não aplicável'' de um
valor imputado, e o valor imputado nunca é zero.

A validação respeita o tempo, com janela expansível: na regressão, seis \textit{folds}
trimestrais (o primeiro treina em 2013T1--T2 e valida em 2013T3; o último valida em
2014T4); na classificação, três \textit{folds} anuais (validação em 2009, 2010 e 2011). O
\textit{pipeline} é ajustado de novo em cada \textit{fold}, apenas com o treino desse
\textit{fold}. No primeiro \textit{fold} da regressão verificou-se que a mediana aprendida
de \texttt{dti} é a do treino (16,98) e não a da validação (16,57), que depois da
padronização a média é nula no treino mas não na validação, e que voltar a ajustar só com
o treino produz exatamente a mesma transformação da validação. Na Parte~2, o
\textit{pipeline} é colocado antes de cada modelo, de forma que o ajuste do
pré-processamento fique sempre dentro do treino de cada \textit{fold}.

O mesmo tipo de \textit{pipeline} foi também implementado em R com o pacote \texttt{recipes}
(\texttt{R/pipeline\_pca.Rmd}), em que \texttt{prep} corresponde ao ajuste e \texttt{bake} à
transformação: logaritmo nas colunas com cauda longa, imputação pela mediana e padronização,
sobre as 14 variáveis numéricas da classificação usadas na análise exploratória. Nos três
\textit{folds} anuais da classificação, as mesmas verificações confirmam que os parâmetros vêm
só do treino (por exemplo, a mediana de \texttt{dti} aprendida no primeiro \textit{fold} é
12,4, a do treino, e não 12,6, a da validação) e que incluir a validação no ajuste alteraria a
transformação. Acrescentando um passo de PCA à mesma receita, o R reproduz exatamente o PCA
feito em Python na análise exploratória (21,6\% da variância na primeira componente e sete
componentes para 80\%), o que serve de verificação cruzada das duas implementações.
